\section{What We Know}
\subsection{Causes of Unfairness}
Even before we precisely specify what we mean by ``fairness'', we can identify common distortions that can lead off-the-shelf machine learning techniques to produce behavior that is intuitively unfair. These include:
\begin{enumerate}
\item \textbf{Bias Encoded in Data}: Often, the training data that we have on hand already includes human biases. For example, in the problem of recidivism prediction used to inform bail and parole decisions, the goal is to predict whether an inmate, if released, will go on to commit another crime within a fixed period of time. But we do not have data on who commits crimes --- we have data on who is arrested. There is reason to believe that arrest data --- especially for drug crimes --- is skewed towards minority populations that are policed at a higher rate \cite{rothwell2014war}. Of course, machine learning techniques are designed to fit the data, and so will naturally replicate any bias already present in the data.  There is no reason to expect them to \emph{remove} existing bias.
\item \textbf{Minimizing Average Error Fits Majority Populations}: Different populations of people have different distributions over features, and those features have different relationships to the label that we are trying to predict. As an example, consider the task of predicting college performance based on high school data. Suppose there is a majority population and a minority population. The majority population employs SAT tutors and takes the exam multiple times, reporting only the highest score. The minority population does not. We should naturally expect both that SAT scores are higher amongst the majority population, and that their relationship to college performance is differently calibrated compared to the minority population. But if we train a group-blind classifier to minimize overall error, if it cannot simultaneously fit both populations optimally, it will fit the majority population. This is because --- simply by virtue of their numbers --- the fit to the majority population is more important to overall error than the fit to the minority population. This leads to a different (and higher) distribution of errors in the minority population. This effect can be quantified, and can be partially alleviated via concerted data gathering efforts \cite{CJS18}.
\item \textbf{The Need to Explore}: In many important problems, including recidivism prediction and drug trials, the data fed into the prediction algorithm depends on the actions that algorithm has taken in the past. We only observe whether an inmate will recidivate \emph{if we release him}. We only observe the efficacy of a drug on patients to whom it is assigned. Learning theory tells us that in order to effectively learn in such scenarios, we need to \emph{explore} --- i.e. sometimes take actions we believe to be sub-optimal in order to gather more data. This leads to at least two distinct ethical questions. First, when are the individual costs of exploration borne disproportionately by a certain sub-population? Second, if in certain (e.g. medical) scenarios, we view it as immoral to take actions we believe to be sub-optimal for any particular patient, how much does this slow learning, and does this lead to other sorts of unfairness?
\end{enumerate}

\subsection{Definitions of Fairness}
With a few exceptions, the vast majority of work to date on fairness in machine learning has focused on the task of batch classification. At a high level, this literature has focused on two main families of definitions\footnote{There is also an emerging line of work that considers \emph{causal} notions of fairness (see e.g., \cite{kilbertus2017avoiding, kusner2017counterfactual, nabi2018fair}). We intentionally avoided discussions of this potentially important direction because it will be the subject of its own CCC visioning workshop.}: \emph{statistical} notions of fairness and \emph{individual} notions of fairness.  We briefly review what is known about these approaches to fairness, their advantages, and their shortcomings.

\subsubsection{Statistical Definitions of Fairness}
Most of the literature on fair classification focuses on \emph{statistical} definitions of fairness. This family of definitions fixes a
small number of protected demographic groups $\mathcal{G}$ (such as racial groups), and then ask for (approximate)
parity of some statistical measure across all of these groups. Popular measures include raw positive classification rate, considered in work such as \cite{CV10,KAS11,awareness,feldman2015certifying} (also sometimes known as \emph{statistical parity} \cite{awareness}), false positive and false negative rates \cite{Chou16,KMR16,HPS16,ZVGG17} (also sometimes known as equalized odds \cite{HPS16}), and positive predictive value \cite{Chou16,KMR16} (closely related to equalized calibration when working with real valued risk scores). There are others --- see e.g. \cite{berksurvey} for a more exhaustive enumeration. This family of fairness definitions is attractive because it is simple, and definitions from this family can be achieved without making any assumptions on the data and can be easily verified. However, statistical definitions of fairness do not on their own give meaningful guarantees to individuals or structured subgroups of the protected demographic groups. Instead they give guarantees to ``average'' members of the protected groups. (See \cite{awareness} for a litany of ways in which statistical parity and similar notions can fail to provide meaningful guarantees, and \cite{gerrymandering} for examples of how some of these weaknesses carry over to definitions which equalize false positive and negative rates.) Different statistical measures of fairness can be at odds with one another. For example, \cite{Chou16} and \cite{KMR16} prove a fundamental impossibility result: except in trivial settings, it is impossible to simultaneously equalize false positive rates, false negative rates, and positive predictive value across protected groups. Learning subject to statistical fairness constraints can also be computationally hard \cite{hardnesspaper}, although practical  algorithms of various sorts are known \cite{HPS16,ZVGG17,MSRpaper}.
\subsubsection{Individual Definitions of Fairness}
Individual notions of fairness, on the other hand, ask for constraints that bind on specific pairs of individuals, rather than on a quantity that is averaged over groups. For example,  \cite{awareness} give a definition which roughly corresponds to the constraint that ``similar individuals should be treated similarly'', where similarity is defined with respect to a task-specific metric that must be determined on a case by case basis. \cite{JKMR16} suggest a definition which roughly corresponds to ``less qualified individuals should not be favored over more qualified individuals'', where quality is defined with respect to the \emph{true} underlying label (unknown to the algorithm). However, although the semantics of these kinds of definitions can be more meaningful than statistical approaches to fairness, the major stumbling block is that they seem to require making significant assumptions. For example, the approach of \cite{awareness} pre-supposes the existence of an agreed upon similarity metric, whose definition would itself seemingly require solving a non-trivial problem in fairness, and the approach of \cite{JKMR16} seems to require strong assumptions on the functional form of the relationship between features and labels in order to be usefully put into practice. These obstacle are serious enough that it remains unclear whether individual notions of fairness can be made practical --- although attempting to bridge this gap is an important and ongoing research agenda.  